\section{Neural Experiment and Statistical Design}
\subsection{Network, Feedback and Calibration}
The FlyWire v783-derived model\cite{dorkenwald2024neuronal} has 138,639 LIF neurons and 15,091,983 signed, contact-weighted connection rows. We use Shiu parameters\cite{shiu2024drosophila} in Brian~2\cite{stimberg2019brian2}, with 0.1-ms steps, 2.2-ms sensory refractoriness and state/weight restoration before each episode. Assumed Gaussian RFs stimulate 104 LC4 and 210 LPLC2 cells at up to 150\,Hz; other external rates are zero. Eight descending cells provide control features for the timed task (Fig.~\ref{fig:v10f-system}(a--c)).

Between events, $\dot v=(v_0-v+g)/(20\,\mathrm{ms})$ and $\dot g=-g/(5\,\mathrm{ms})$. Threshold $-45$\,mV resets $v$ to $v_0=-52$\,mV and $g$ to zero; both equations pause during refractoriness. Recurrent events add signed contact count times 0.275\,mV to $g$ after 1.8\,ms. External events add 68.75\,mV to voltage. These are functional model parameters, not measured device-fault characteristics.

Input begins at 100\,ms; sensing and feedback update every 20\,ms. Bearing is world-relative angle minus heading; looming is the change of $\arctan(0.05/d)$ at actual separation $d$. The pursuer starts 0.25\,m away at $\pm30^\circ$ and moves at 0.3\,m/s. Each motion segment aims it at the agent's segment-start position. Both positions advance linearly using the agent's starting heading; heading updates at segment end, capped at $\pm100^\circ$. A frame-end decision affects the next frame. The first trigger schedules a 0.35-m/s, heading-fixed burst 250\,ms later for 450\,ms; otherwise speed is 0.05\,m/s. Segments split at exact actuator events; first contact with radius 0.05\,m is solved within each segment. Failure latches, while traces continue to $H=1$\,s for resource accounting. Post-contact activity is counterfactual and excluded from warnings.

For each map, 20 passive development episodes (four seeds, five bearings from $-60^\circ$ to $60^\circ$) fit an affine steering score $z$. Standardized five-frame means of eight command rates over 200--400\,ms predict bearing/$60^\circ$, using ridge penalty 10 and an unpenalized intercept. The trigger threshold is half the tenth percentile of development episode mean command rates over 200--300\,ms, floored at 2.5\,Hz. Before the cue or below threshold, steering is zero. Otherwise proportional/relay laws apply $-600\operatorname{clip}(z,-1,1)$/$-600\operatorname{sgn}(z)$ deg/s outside the burst. Healthy held-out validation RMSE is $65.9^\circ/67.8^\circ/67.2^\circ$ for RF-0/A/B, over eight seed clusters each: $z$ is not an accurate closed-loop bearing estimate.

\subsection{Pairing, Cohorts and Endpoints}
Shared uniforms give nested deletions $w_j^{(f)}=w_j\mathbf1[V_{sj}\geq f]$ and input events $\mathbf1[U_{sk\tau i}<\lambda_{ki}\Delta t]$. Indices denote seed $s$, frame $k$, neural step $\tau$, cell $i$ and connection row $j$; $\Delta t=10^{-4}$\,s. Feedback changes $\lambda$, while input variates and row masks remain paired. Seeds and hashes are traceable, without cross-platform bit-identity guarantees.

RF-0 uses 24 test seed clusters; RF-A/B each use 12 new clusters. Every cluster contains both bearings, two laws and healthy/25\%/50\% deletion: 288/144/144 episodes. New maps independently redraw all 314 azimuths from the same hemisphere prior (left $10$--$90^\circ$, right $-90$--$-10^\circ$), preserving identities, vertical mapping and gradients. Each separately fits the frozen development recipe, uses 32 healthy validation episodes in eight further clusters, then freezes thresholds before testing. All maps, seeds and failures are retained; no test tuning or efficacy stopping occurs. This is transfer with recalibration, not zero-shot transfer or an RF-isolated causal experiment.

With $Y=1$ meaning no contact through $H$, registered primary endpoints are paired loss $L=Y_H(1-Y_F)$, gain $G=(1-Y_H)Y_F$ and activity exposure. Net loss $D=L-G$ is derived. The relay-minus-proportional contrast and warnings are prespecified secondary; trigger-timing and clearance/heading diagnostics are post-hoc. For cumulative spikes $S$ and $T_*=\min(T_F,T_H)$, own-stop and common-duration ratios are
\begin{equation}
R_{\rm stop}=\frac{S_F(T_F)}{S_H(T_H)},\qquad
R_{\rm common}=\frac{S_F(T_*)}{S_H(T_*)}.
\end{equation}
Here $T$ is contact-frame end or $H$, at 20-ms resolution. Common duration remains outcome-dependent and does not equate sensory histories; fixed 400-ms/1-s windows are separate endpoints.

Estimates average both bearings within seed. Pointwise 95\% CIs use 20,000 whole-seed percentile resamples, conditional on stored map/calibration/thresholds (RNG 20261022 for RF-0; 20261023 for transfer). Symmetric-null sign-flip sensitivity tests are Holm-adjusted over two baseline doses, separately from four new map/dose contrasts. Maps and separate angle/action cohorts are not pooled.

Warnings use only data through 400\,ms: unsmoothed 20-ms command rates averaged over 100--400\,ms, or the absolute log spike-count ratio to a law/bearing-specific healthy reference. Each RF's healthy validation separately sets fifth-percentile command/forecast and 95th-percentile two-sided activity cutoffs, frozen before its test cohort; an upper-only comparator flags activity above twice its healthy reference. A held-command forecast uses exact simulated geometry and pending actuation, retaining the current turn; its schedule-only baseline sets future turn to zero. These information-rich forecasts are not neural-only detectors.
